%% kapa: exam=1 round=21 year=2010 period=1 form=A subject=civil_law subject_ko=민법 seq=4 number=4
%% source: Q-Net 감정평가사 기출문제 (cst00309) artlSeq=5024181 file=extracted/1008/5024181/2010년 제21회 감정평가사 1차시험 기출문제/1교시.hwp
%% conversion: kordoc 4.24.51 -> tools/gen1.py (auto); review=proofread-sonnet; flags=none
\begin{question}{4}
등기청구권에 관한 설명으로 옳지 않은 것은? (다툼이 있으면 판례에 의함)
\begin{choices}
  \item 등기권리자가 등기의무자에게 등기신청에 협력하여 줄 것을 요구할 수 있는 실체법상의 권리이다.
  \item 매매로 인한 소유권이전등기청구권의 양도는 특별한 사정이 없는 한 채무자의 동의나 승낙을 받아야 대항할 수 있다.
  \item 부동산 매수인이 부동산을 인도받아 사용ㆍ수익하다가 제3자에게 그 부동산을 처분하고 점유를 승계하여 준 경우, 소유권이전등기청구권의 소멸시효는 진행한다.
  \item 가등기에 기한 소유권이전등기청구권이 시효완성으로 소멸되었다면, 가등기 이후에 그 부동산을 취득한 제3자는 가등기권자에 대하여 그 가등기의 말소를 구할 수 있다.
  \item 시효기간 진행 중 소유명의자가 변경되고 그 후에 시효가 완성되면 시효완성자는 시효완성 당시의 소유자를 상대로 등기청구권을 행사할 수 있다.
\end{choices}
\end{question}
\answer{3}
